% Part: set-theory
% Chapter: spine
% Section: idea

\documentclass[../../../include/open-logic-section]{subfiles}

\begin{document}

\olfileid{sth}{spine}{valpha}

\olsection{পর্যায়গুলিকে $V_\alpha$-এর মাধ্যমে সংজ্ঞায়িত করা}

\olref[sth][ordinals][]{chap}-এ আমরা সুক্রমবিন্যাস
ও (ফন নয়মানের) ক্রমসংখ্যা সংজ্ঞায়িত করেছি। এবার
সেগুলির সাহায্যে সেটের স্তরক্রমকে \emph{নিজস্ব
তত্ত্বের মধ্যেই} চিহ্নিত করব। এ জন্য মনে করি,
\olref[sth][ordinals][opps]{sec}-এ উত্তরসূরি ও
সীমা ক্রমসংখ্যার ধারণা সংজ্ঞায়িত হয়েছে। সেই
ধারণাগুলি ব্যবহার করে নিচের সংজ্ঞা দিই:

\begin{defn}\ollabel{defValphas}
	\begin{align*}
	V_\emptyset &\defis \emptyset\\
	V_{\ordsucc{\alpha}} &\defis \Pow{V_\alpha} & &
	\text{যেকোনো ক্রমসংখ্যা }\alpha\text{-এর জন্য}\\
	V_{\alpha} &\defis \bigcup_{\gamma < \alpha} V_\gamma & &
	\text{যখন }\alpha\text{ সীমা ক্রমসংখ্যা}
\end{align*}
\end{defn}
\noindent
এটি ক্রমসংখ্যার উপর \emph{ট্রান্সফাইনাইট পুনরাবৃত্তি}
দিয়ে সংজ্ঞা। স্বাভাবিক সংখ্যার উপর পুনরাবৃত্তি
দিয়ে অপেক্ষক সংজ্ঞায়নের সঙ্গে এর তুলনা করা উচিত।\footnote{যোগ,
গুণ ও ঘাতের সংজ্ঞাগুলি দেখুন
\olref[sfr][infinite][dedekind]{sec}-এ।} স্বাভাবিক
সংখ্যার মতো এখানে ভিত্তি ও উত্তরসূরি ক্ষেত্র
সংজ্ঞায়িত করি; কিন্তু ক্রমসংখ্যার ক্ষেত্রে
\emph{সীমা} ক্ষেত্রের আচরণও বলতে হয়।

$V_\alpha$-গুলির এই সংজ্ঞা একটি গুরুত্বপূর্ণ
মাইলফলক। এত দিন সেটকে \emph{পর্যায়ে পর্যায়ে}
গঠনের ধারণা দিয়ে আমরা স্তরক্রমটি অনানুষ্ঠানিকভাবে
সমর্থন করেছি। $V_\alpha$-গুলি কার্যত সেই পর্যায়গুলিই।
তবে গুরুত্বপূর্ণ হলো, এটি পর্যায়গুলির একটি
\emph{তত্ত্ব-অভ্যন্তরীণ} চরিত্রায়ন। তত্ত্বের কোনো
সম্ভাব্য \emph{মডেল} প্রস্তাব করার বদলে আমরা
আমাদের সেটতত্ত্বের \emph{ভিতরেই} পর্যায়গুলি
সংজ্ঞায়িত করব।

\end{document}
